inde ut in art. 2. concludimus, $Y$ cum $y$ identicam esse. quoties $y$ quoque sit
functio algebraica integra ordinem $n$ non transscendens, aut saltem ipsius $y$ vice
fungi posse, si $y$ in seriem secundum potestates ipsius $t$ progredientem conversa
tantam convergentiam exhibeat, ut terminos altiorum ordinum negligere liceat.

\SeventhArticle{8}
Iam ad eruendum integrale $\int Y\,\mathrm{d}t$ singulas partes ipsius $Y$ consideremus.
Designemus productum,
\[
(t-a)(t-a')(t-a'')(t-a''')\ldotp\ldotp\ldotp\ldotp(t-a^{(n)})
\]
per $T$, fiatque per evolutionem huius producti
\[
T=t^{n+1}+\alpha t^n+\alpha't^{n-1}+\alpha''t^{n-2}+\text{etc.}+\alpha^{(n)}
\]
Numerator fractionis, per quam, in parte prima ipsius $Y$, multiplicata est $A$, fit
$=\frac{T}{t-a}$; numeratores in partibus sequentibus perinde sunt
$\frac{T}{t-a'},\frac{T}{t-a''},\frac{T}{t-a'''}$ etc.
Denominatores vero nihil aliud sunt, nisi valores determinati horum numeratorum,
si resp. statuitur $t=a$, $t=a'$, $t=a''$, $t=a'''$ etc.: denotemus hos denominatores
resp. per $M,M',M'',M'''$ etc., ita ut sit
\[
\begin{aligned}&Y=\frac{AT}{M(t-a)}+\frac{A'T}{M'(t-a')}+\frac{A''T}{M''(t-a'')}\\ &+\text{etc.}+\frac{A^{(n)}T}{M^{(n)}(t-a^{(n)})}\end{aligned}
\]
Quum fiat $T=0$, pro $t=a$, habemus aequationem identicam
\[
a^{n+1}+\alpha a^n+\alpha'a^{n-1}+\alpha''a^{n-2}+\text{etc.}+\alpha^{(n)}=0
\]
adeoque
\[
\begin{aligned}
T={}&t^{n+1}-a^{n+1}+\alpha(t^n-a^n)+\alpha'(t^{n-1}-a^{n-1})
\\ &+\alpha''(t^{n-2}-a^{n-2})+\text{etc.}\\
&+\alpha^{(n-1)}(t-a)
\end{aligned}
\]
Hinc dividendo per $t-a$ fit
\[
\begin{aligned}
\frac{T}{t-a}={}&t^n+at^{n-1}+aat^{n-2}+a^3t^{n-3}+\text{etc.}+a^n\\[4pt]
&+\alpha t^{n-1}+\alpha at^{(n-2)}+\alpha aat^{n-3}+\text{etc.}+\alpha a^{n-1}\\[4pt]
&\hspace{1.01em}+\alpha't^{n-2}+\alpha'at^{n-3}+\text{etc.}+\alpha'a^{n-2}\\[4pt]
&\hspace{2.01em}+\alpha''t^{n-3}+\text{etc.}+\alpha''a^{n-3}\\[4pt]
&\hspace{3.01em}+\text{etc. etc.}\\[4pt]
&\hspace{4.01em}+\alpha^{(n-1)}
\end{aligned}
\]
\SourceLeaf{183}
Valor huius functionis pro $t=a$ colligitur
\[
\begin{aligned}&=(n+1)a^n+n\alpha a^{n-1}+(n-1)\alpha'a^{n-2}\\ &+(n-2)\alpha''a^{n-3}
+\text{etc.}+\alpha^{(n-1)}\end{aligned}
\]
Hinc $M$ aequalis valori ipsius $\frac{\mathrm{d}T}{\mathrm{d}t}$ pro $t=a$, uti etiam
aliunde constat. Perinde $M',M'',M'''$, etc. erunt valores ipsius
$\frac{\mathrm{d}T}{\mathrm{d}t}$ pro $t=a'$, $t=a''$, $t=a'''$ etc.

Porro invenimus valorem integralis $\int\frac{T\,\mathrm{d}t}{t-a}$, a $t=0$ usque ad $t=1$,
\[
\begin{aligned}
={}&\frac{1}{n+1}+\frac{a}{n}+\frac{aa}{n-1}+\frac{a^3}{n-2}+\text{etc.}+a^n\\[4pt]
&+\frac{\alpha}{n}+\frac{\alpha a}{n-1}+\frac{\alpha aa}{n-2}+\text{etc.}+\alpha a^{n-1}\\[4pt]
&\hspace{5.01em}+\frac{\alpha'}{n-1}+\frac{\alpha'a}{n-2}+\text{etc.}+\alpha'a^{n-2}\\[4pt]
&\hspace{6.01em}+\frac{\alpha''}{n-2}+\text{etc.}+\alpha''a^{n-3}\\[4pt]
&\hspace{7.01em}+\text{etc. etc.}\\[4pt]
&\hspace{8.01em}+\alpha^{(n-1)}
\end{aligned}
\]
quos terminos ordine sequenti disponemus:
\[
\begin{aligned}
&a^n+\alpha a^{n-1}+\alpha'a^{n-2}+\alpha''a^{n-3}+\text{etc.}+\alpha^{(n-1)}\\[4pt]
&\quad+\tfrac12(a^{n-1}+\alpha a^{n-2}+\alpha'a^{n-3}+\text{etc.}+\alpha^{(n-2)})\\[4pt]
&\quad+\tfrac13(a^{n-2}+\alpha a^{n-3}+\alpha'a^{n-4}+\text{etc.}+\alpha^{(n-3)})\\[4pt]
&\quad+\tfrac14(a^{n-3}+\alpha a^{n-4}+\alpha'a^{n-5}+\text{etc.}+\alpha^{(n-4)})\\[4pt]
&\quad+\text{etc.}\\[4pt]
&\quad+\frac{1}{n-1}(aa+\alpha a+\alpha')\\[4pt]
&\quad+\frac{1}{n}(a+\alpha)\\[4pt]
&\quad+\frac{1}{n+1}
\end{aligned}
\]
Sed manifesto eadem quantitas prodit, si in producto e multiplicatione functionis
$T$ in seriem infinitam
\[
t^{-1}+\tfrac12t^{-2}+\tfrac13t^{-3}+\tfrac14t^{-4}\ \text{etc.}
\]
orto, reiectis omnibus terminis, qui implicant potestates ipsius $t$ exponentibus
negativis (sive brevius, in producti parte ea, quae est functio integra ipsius $t$)
pro $t$ scribitur $a$. Supponamus itaque. fieri\HistoricalNote{19}
\[
T(t^{-1}+\tfrac12t^{-2}+\tfrac13t^{-3}+\tfrac14t^{-4}+\text{etc.})=T'+T''
\]
\SourceLeaf{184}
ita ut $T'$ sit functio integra ipsius $t$ in hoc producto contenta, $T''$ vero pars
altera, scilicet series descendens in infinitumque excurrens. Quo facto valor integralis
$\int\frac{T\,\mathrm{d}t}{t-a}$ a $t=0$ usque ad $t=1$ aequalis erit valori functionis $T'$
pro $t=a$. Quodsi itaque valores determinatos functionis
\[
\frac{T'}{\left(\dfrac{\mathrm{d}T}{\mathrm{d}t}\right)}
\]
pro $t=a$, $t=a'$, $t=a''$, $t=a'''$ etc. usque ad $t=a^{(n)}$ resp. per
$R,R',R'',R'''\ldotp\ldotp\ldotp\ldotp R^{(n)}$ denotamus, integrale $\int Y\,\mathrm{d}t$ a $t=0$ usque ad $t=1$ fiet
\[
=RA+R'A'+R''A''+\text{etc.}+R^{(n)}A^{(n)}
\]
quod per $\Delta$ multiplicatum exhibebit valorem vel verum vel approximatum integralis
$\int y\,\mathrm{d}x$ ab $x=g$ usque ad $x=g+\Delta$.

\SeventhArticle{9}
Hae operationes aliquanto facilius perficiuntur, si loco indeterminatae $t$
introducitur alia $u=2t-1$. Scribimus quoque brevitatis caussa $b=2a-1$,
$b'=2a'-1$, $b''=2a''-1$ etc. Transeat $T$, substituto pro $t$ valore
$\tfrac12u+\tfrac12$, in $\frac{U}{2^{n+1}}$, sive sit
\[
U=(u-b)(u-b')(u-b'')\ldotp\ldotp\ldotp\ldotp(u-b^n)
\]
Erit itaque $\frac{\mathrm{d}T}{\mathrm{d}t}=\frac{1}{2^n}\cdot\frac{\mathrm{d}U}{\mathrm{d}u}$,
adeoque $M,M',M''$ etc. valores determinati ipsius
$\frac{1}{2^n}\cdot\frac{\mathrm{d}U}{\mathrm{d}u}$, si deinceps statuitur $u=b$, $u=b'$, $u=b''$ etc.

Quum series $t^{-1}+\tfrac12t^{-2}+\tfrac13t^{-3}+\tfrac14t^{-4}+\text{etc.}$ nihil aliud sit quam
$\log.\frac{1}{1-t^{-1}}=\log\frac{1+u^{-1}}{1-u^{-1}}$:
per substitutionem $t=\tfrac12u+\tfrac12$ necessario transibit in
$2u^{-1}+\tfrac23u^{-3}+\tfrac25u^{-5}+\tfrac27u^{-7}+\text{etc.}$ Quodsi itaque statuimus
\[
U(u^{-1}+\tfrac13u^{-3}+\tfrac15u^{-5}+\tfrac17u^{-7}+\text{etc.})=U'+U''
\]
ita ut $U'$ sit functio integra ipsius $u$ in hoc producto contenta, $U''$ vero pars
altera, puta series descendens infinita, patet esse
\[
T'+T''=\frac{1}{2^n}(U'+U'')
\]
Sed manifesto $T'$, tamquam functio integra ipsius $t$, per substitutionem
$t=\tfrac12u+\tfrac12$ necessario functionem integram ipsius $u$ producet: contra $T''$,
quae non continet nisi potestates negativas ipsius $t$, per eandem substitutionem
tantummodo potesta\SourceLeaf{185}tes negativas ipsius $u$ gignet. Quam ob rem $U'$ nihil aliud erit quam $2^nT'$ per hanc substitutionem transformata, ac perinde $U''$ producta erit ex $2^nT''$. Nihil itaque intererit, sive in $\dfrac{T'}{\left(\dfrac{\mathrm{d}T}{\mathrm{d}t}\right)}$ substituamus $t=a$, sive in $\dfrac{U'}{\left(\dfrac{\mathrm{d}U}{\mathrm{d}u}\right)}$ faciamus $u=b$, unde colligimus, $R, R', R'', R'''$ etc. etiam esse valores determinatos functionis $\dfrac{U'}{\left(\dfrac{\mathrm{d}U}{\mathrm{d}u}\right)}$ pro $u=b$. $u=b',\ u=b'',\ u=b'''$ etc.

\SeventhArticle{10}
Antequam ulterius progrediamur, haecce praecepta per exemplum illustrabimus. Sit $n=5$, statuamusque $a=0,\ a'=\frac15,\ a''=\frac25,\ a'''=\frac35,\ a''''=\frac45,\ a'''''=1$. Hinc fit
\[
T=t^6-3t^5+\tfrac{17}{5}t^4-\tfrac95t^3+\tfrac{274}{625}tt-\tfrac{24}{625}t
\]
Multiplicando per $t^{-1}+\frac12t^{-2}+\frac13t^{-3}+\frac14t^{-4}+$ etc. obtinemus
\[
T'=t^5-\tfrac52t^4+\tfrac{67}{30}t^3-\tfrac{17}{20}tt+\tfrac{913}{7500}t-\tfrac{19}{7500}
\]
Valores itaque co\"efficientium $R, R', R'', R''', R'''', R'''''$ exprimuntur per functionem fractam
\[
\frac{\begin{gathered}t^5-\frac25t^4+\frac{67}{30}t^3\\-\frac{17}{20}tt+\frac{913}{7500}t-\frac{19}{7500}\end{gathered}}
{\begin{gathered}6t^5-15t^4+\frac{68}{5}t^3\\-\frac{27}{5}tt+\frac{548}{625}t-\frac{24}{625}\end{gathered}}
\]
in qua pro $t$ deinceps substituendi sunt valores $0,\ \frac15,\ \frac25,\ \frac35,\ \frac45,\ 1$. Aliquanto brevior est methodus altera, quae suppeditat $b=-1,\ b'=-\frac35,\ b''=-\frac15,\ b'''=\frac15,\ b''''=\frac35,\ b'''''=1$
\[
\begin{aligned}
U&=u^6-\tfrac75u^4+\tfrac{259}{625}uu-\tfrac9{625}\\
U'&=u^5-\tfrac{16}{15}u^3+\tfrac{277}{1875}u
\end{aligned}
\]
unde $R.\ R',\ R''$ etc. erunt valores functionis fractae
\[
\frac{u^4-\frac{16}{15}uu+\frac{277}{1875}}
{6u^4-\frac{28}{5}uu+\frac{518}{625}}
\]
pro $u=-1,\ u=-\frac35,\ u=-\frac15$ etc. Utraque methodus eosdem numeros profert, quos in art. 4. ex Harmonia Mensurarum tradidimus. Ceterum in casu tali, qualem hocce exemplum sistit, ubi $a, a', a''$ etc. sunt quantitates rationales, valores denominatoris $\frac{\mathrm{d}T}{\mathrm{d}t}$ commodius in forma primitiva computantur, puta $(a-a')(a-a'').(a-a''')\ldotp\ldotp\ldotp\ldotp(a-a^{(n)})$ pro $t=a$ ac perinde de reliquis. Idem valet de denominatore $\frac{\mathrm{d}U}{\mathrm{d}u}$, qui pro $u=b$ fit $=(b-b')(b-b'')(b-b''')\ldotp\ldotp\ldotp\ldotp(b-b^{(n)})$.
